\SourceLeaf{358}%
zen Erdkörpers nur kleine Grössen sein können, solche Umwandlungen der Formeln, welche die Geschmeidigkeit und Bequemlichkeit derselben sehr vergrössern; ja, wenn gleich diese Umwandlungen eigentlich nur Näherungsformeln sind, so können sie doch nicht bloss eben so grosse, sondern selbst grössere Schärfe gewähren, als die absolut strengen Formeln, was man. nicht paradox finden wird, wenn man erwägt, dass die letztern doch immer vermittelst der trigonometrischen Tafeln zur Ausübung kommen müssen, deren Schärfe keine absolute, sondern durch die Anzahl der Decimalzifern begrenzt ist. Unter den verschiedenen in der ersten Abhandlung mitgetheilten für den angedeuteten Zweck bestimmten Formeln zeichnet sich nun besonders die am Schluss derselben aufgeführte Combination dadurch aus, dass sie den Zusammenhang jener sechs Quantitäten in der zur Rechnung möglich bequemsten Gestalt aufstellt, und eine Schärfe gewährt, die auch bei den grössten wirklich messbaren Dreiecken überflüssig ausreicht. Es musste dadurch das Verlangen nach dem Besitz analoger unmittelbar für die Ellipsoidfläche geltender Formeln erweckt werden, und die Entwicklung derselben bildet den Hauptinhalt der gegenwärtigen \textit{zweiten} Abhandlung.

Während die Auffindung der erwähnten für die Kugelfläche gültigen Formeln auf ganz elementarischen Sätzen beruhete, erfordert hingegen die Ermittlung ihrer Gegenstücke auf der Ellipsoidfläche eine Reihe ziemlich verwickelter Operationen, und es muss daher ohne Zweifel angenehm sein, wenn mehr als Ein Weg zu demselben Ziele zu gelangen nachgewiesen wird. Der Verf., welcher alle diese Untersuchungen schon vor mehr als dreissig Jahren zu seinem Privatgebrauch durchgeführt, und nur bisher zur Veröffentlichung noch keine besondere Veranlassung gefunden hatte, theilt nun in der vorliegenden Abhandlung \textit{zwei} unter sich durchaus verschiedene, aber zuletzt zu ganz gleichen Resultaten führende Ableitungsarten mit, von denen eine in der Theorie der conformen Übertragung der Ellipsoidfläche auf die Kugelfläche wurzelt. In dieser Beziehung schliesst sich die zweite Abhandlung auch an die erste an, obwohl übrigens beide insofern als gänzlich unabhängig von einander zu betrachten sind, als man freie Wahl behält, die geodaetischen Rechnungen entweder bloss nach der in der ersten Abhandlung, oder bloss nach der in der zweiten Abhandlung gelehrten Methode zu führen.

Die Aufgabe, von der geographischen Lage eines Punkts auf der Sphäroidfläche zu der eines andern Punktes überzugehen, der mit jenem durch eine geo%
